% -*- Mode: latex -*- %

\section{Proofs}

\subsection{Proof of Theorem \ref{theorem:normality}}
\label{sec:proof-normality}

We formally state the smoothness condition needed for Theorem \ref{theorem:normality}.

\begin{definition}[Smooth enough losses]
  \label{definition:smooth-enough}
  The loss $\loss_\theta$ is \emph{smooth enough} if
  \begin{enumerate}[label=(\roman*),leftmargin=1.5em]
  \item  
  \label{item:loss-differentiability} the losses
    $\loss_\theta(x, y)$ and $\loss_\theta(x, f(x))$ are
    differentiable at $\theta\opt$ for $\P$-almost
    every $(x, y)$,
  \item \label{item:locally-lipschitz}
    the losses $\loss_\theta$ are
    locally Lipschitz around $\theta\opt$: there is a neighborhood
    of $\theta\opt$ such that
    $\loss_\theta(x, y)$ is $\lipobj(x, y)$-Lipschitz
    and $\loss_{\theta}(x, f(x))$ is $\lipobj(x)$-Lipschitz in $\theta$,
    where $\E[\lipobj(X, Y)^2 + \lipobj(X)^2] < \infty$,
  \item the population losses
    $L(\theta) = \E[\loss_\theta(X, Y)]$ and
    $L^f(\theta) = \E[\loss_\theta(X, f(X))]$ have Hessians $H_{\theta\opt} = \nabla^2 L(\theta\opt) \succ 0$ and
    $H_{\theta\opt}^f = \nabla^2 L^f(\theta\opt)$.
  \end{enumerate}
\end{definition}


The proof follows that of \citet[Theorem
  5.23]{VanDerVaart98}, with the modifications necessary
to address the different observation model and empirical objective.  Throughout the proof, we recall the labeling
conventions that a tilde $\wt{\cdot}$ indicates the unlabeled sample
$\{\wt{X}_i\}_{i = 1}^N$ and a superscripted $f$ (i.e.\ $\cdot^f$) indicates
using the machine-learning predictions $f$.

With this notational setting,
given a function $g$, we use the shorthand notation
\begin{align}
  \E_n g &:= \frac{1}{n} \sum_{i=1}^n g(X_i,Y_i), \quad \GG_n g = \sqrt{n}\left(\E_n g - \E[g(X,Y)]\right); \label{eq:Gn}\\
  \widetilde \E_N^f g &:= \frac{1}{N} \sum_{i=1}^N g(\widetilde X_i, f(\widetilde X_i)),\quad \widetilde \GG_N^f g := \sqrt{N} ( \widetilde \E_N^f g - \E[g(X,f(X))]); \label{eq:GNf}\\
  \E_n^f g &:= \frac{1}{n} \sum_{i=1}^n g( X_i, f(X_i)), \quad \GG_n^f g := \sqrt{n}  (\E_n^f g - \E[g(X,f(X))]) \label{eq:Gnf}.
\end{align}

The differentiability of $\loss_\theta(x,y)$ and
$\loss_{\theta}(x, f(x))$ at $\theta\opt$ (Definition~\ref{definition:smooth-enough}\ref{item:loss-differentiability}) and the local Lipschitzness of the losses
(Definition~\ref{definition:smooth-enough}\ref{item:locally-lipschitz})
imply that
for every (possibly random) sequence $h_n = O_P(1)$, we have
\begin{equation*}
  \GG\left[\sqrt{n} \left(\loss_{\theta\opt + \frac{h_n}{\sqrt{n}}} - \loss_{\theta\opt} \right) - h_n^\top \nabla \loss_{\theta\opt} \right]
  \cp 0
\end{equation*}
for any of the empirical processes $\GG\in \{\GG_n, \widetilde \GG_N^f,
\GG_n^f \}$ (see \cite[Lemma 19.31]{VanDerVaart98}).  By applying a second-order Taylor expansion, this implies
\begin{align*}
  n \E_n\left(\loss_{\theta\opt + \frac{h_n}{\sqrt{n}}} - \loss_{\theta\opt}\right)
  & =
  n \left(L\Big(\theta\opt + \frac{h_n}{\sqrt{n}}\Big)
  - L(\theta\opt)\right)
  + h_n^\top \GG_n \nabla \loss_{\theta\opt} + o_P(1) \\
  & = \half h_n^\top H_{\theta\opt} h_n
  + h_n^\top \GG_n \nabla\loss_{\theta\opt} + o_P(1).
\end{align*}
Similarly, we have the two equalities
\begin{align*}
  \hat\lambda n \widetilde \E_N^f\left(\loss_{\theta\opt + \frac{h_n}{\sqrt{n}}} - \loss_{\theta\opt}\right) &= (\lambda + o_P(1))\left(\frac{1}{2} h_n^\top H_{\theta\opt}^f h_n + \sqrt{\frac{n}{N}} h_n ^\top \widetilde \GG_N^f \nabla \loss_{\theta\opt} + o_P(1)\right)\\
  & = \lambda \left(\frac{1}{2} h_n^\top H_{\theta\opt}^f h_n + \sqrt{\frac{n}{N}} h_n ^\top \widetilde \GG_N^f \nabla \loss_{\theta\opt}\right) + o_P(1),
  ~~ \mbox{and} \\
 - \hat\lambda n \E_n^f\left(\loss_{\theta\opt + \frac{h_n}{\sqrt{n}}} - \loss_{\theta\opt}\right) &= (\lambda + o_P(1))\left( -\frac{1}{2} h_n^\top H_{\theta\opt}^f h_n - h_n ^\top \GG_n^f  \nabla \loss_{\theta\opt} + o_P(1) \right)\\
 &= \lambda \left( -\frac{1}{2} h_n^\top H_{\theta\opt}^f h_n - h_n ^\top \GG_n^f  \nabla \loss_{\theta\opt}\right) + o_P(1)  \label{eq:G_nf}.
\end{align*}
%Note that for the second equation we multiplied the expression by $\sqrt{\frac{n}{N}}$.

Adding the preceding three equations and recalling the definition
that 
\[\LPP_\lambda(\theta) = \E_n \loss_\theta
+ \lambda (\wt{\E}_N^f \loss_\theta - \E_n^f \loss_\theta),\]
we have
\begin{equation*}
  n \left(\LPP_{\hat\lambda}\left(\theta\opt + \frac{h_n}{\sqrt{n}}\right)
  - \LPP_{\hat\lambda}(\theta\opt)\right)
  = \half h_n^\top H_{\theta\opt} h_n
  + h_n^\top \left(\GG_n \nabla\loss_{\theta\opt}
  + \lambda\sqrt{\frac{n}{N}} \widetilde \GG_N^f \nabla\loss_{\theta\opt}
  - \lambda \GG_n^f \nabla\loss_{\theta\opt}\right) + o_P(1).
\end{equation*}

We now evaluate this expression for two particular choices of the sequence
$h_n$. First, consider $h_n\opt = \sqrt{n}(\thetaPP_{\hat\lambda} -
\theta\opt)$.  Because the losses are smooth enough
(Definition~\ref{definition:smooth-enough}), $\theta \mapsto \lambda
\loss_\theta$ is locally Lipschitz uniformly for $\lambda$ in (any) compact set,
and so the consistency $\thetaPP_{\hat\lambda}\to \theta\opt$ implies that
$h_n\opt = O_P(1)$ (\citet[Corollary 5.53]{VanDerVaart98}).
This yields
\begin{equation*}
  n \left(\LPP_{\hat\lambda}(\thetaPP_{\hat\lambda}) - \LPP_{\hat\lambda}(\theta\opt)\right) = \frac{1}{2} {h_n\opt}^\top H_{\theta\opt} h_n\opt + {h_n\opt} ^\top \left(\GG_n \nabla\loss_{\theta\opt} + \lambda\sqrt{\frac{n}{N}} \widetilde \GG_N^f \nabla\loss_{\theta\opt} - \lambda \GG_n^f \nabla\loss_{\theta\opt}\right) + o_P(1).
\end{equation*}
We can perform a similar expansion with
$h_n = -H_{\theta\opt}^{-1}
(\GG_n \nabla\loss_{\theta\opt} + \lambda \sqrt{\frac{n}{N}} \widetilde
\GG_N^f \nabla\loss_{\theta\opt} - \lambda \GG_n^f
\nabla\loss_{\theta\opt})$, which is certainly $O_P(1)$,
yielding
\begin{equation*}
  n \left(\LPP_{\hat\lambda}(\theta\opt -h_n/\sqrt{n}) - \LPP_{\hat\lambda}(\theta\opt)\right)
  = -\half h_n^\top H_{\theta\opt} h_n + o_P(1).
%% &= -\half \left(\GG_n \nabla\loss_{\theta\opt} + \lambda\sqrt{\frac{n}{N}} \widetilde \GG_N^f \nabla\loss_{\theta\opt} - \lambda \GG_n^f \nabla\loss_{\theta\opt}\right)^\top  H_{\theta\opt}^{-1} \left(\GG_n \nabla\loss_{\theta\opt} + \lambda \sqrt{\frac{n}{N}} \widetilde \GG_N^f \nabla\loss_{\theta\opt} - \lambda \GG_n^f \nabla\loss_{\theta\opt}\right) + o_P(1).
\end{equation*}
By the definition of $\thetaPP_{\hat\lambda}$, the left-hand side of the
first equation is smaller than the left-hand side of the second equation,
and so
\begin{align*}
  \half {h_n\opt}^\top H_{\theta\opt} h_n\opt
  - {h_n\opt}^\top H_{\theta\opt} h_n
  \le - \half h_n^\top H_{\theta\opt} h_n + o_P(1),
\end{align*}
or, rearranging,
\begin{equation*}
  \half (h_n\opt - h_n)^\top H_{\theta\opt} (h_n\opt - h_n)
  = o_P(1).
\end{equation*}
As $H_{\theta\opt}$ is positive definite, we thus must have
$h_n\opt = h_n + o_P(1)$, that is,
\begin{equation*}
  \sqrt{n} \left(\thetaPP_{\hat\lambda} - \theta\opt\right)
  = -H_{\theta\opt}^{-1} \left(\GG_n \nabla \loss_{\theta\opt}
  + \lambda \sqrt{\frac{n}{N}} \wt{\GG}_N^f \nabla \loss_{\theta\opt}
  - \lambda \GG_n^f \nabla \loss_{\theta\opt}\right) + o_P(1).
\end{equation*}

It remains to demonstrate the asymptotic normality of the empirical
processes on the right-hand side. This follows by the central limit theorem:
\begin{align*}
&\GG_n \nabla\loss_{\theta\opt} + \lambda \sqrt{\frac{n}{N}} \widetilde \GG_N^f \nabla\loss_{\theta\opt} - \lambda \GG_n^f \nabla\loss_{\theta\opt}\\
&\quad = \lambda \sqrt{\frac{n}{N}} \cdot \frac{1}{\sqrt{N}} \sum_{i=1}^N (\nabla \tilde \loss_{\theta\opt,i}^f - \E[ \nabla \loss_{\theta\opt}^f])  + \frac{1}{\sqrt n} \sum_{i=1}^n \left( \nabla \loss_{\theta\opt,i} - \lambda \nabla \loss_{\theta\opt,i}^f - \E[\nabla \loss_{\theta\opt} -  \lambda\nabla \loss_{\theta\opt}^f]\right)\\
&\quad \cd \normal \left(0, \lambda^2 r \cdot \cov(\nabla \loss_{\theta\opt}^f) + \cov(\nabla \loss_{\theta\opt} - \lambda\nabla \loss_{\theta\opt}^f)\right).
\end{align*}
Thus $H_{\theta\opt}^{-1} \left(\GG_n \nabla\loss_{\theta\opt} + \lambda
\sqrt{\frac{n}{N}} \widetilde \GG_N^f \nabla\loss_{\theta\opt} - \lambda
\GG_n^f \nabla\loss_{\theta\opt}\right)\cd \normal(0,\Sigma^\lambda)$,
completing the proof.

\subsection{Proof of Proposition \ref{prop:consistency-convex}}
\label{sec:proof-consistency-convex}

We begin by stating an auxiliary lemma.

\begin{lemma}
  \label{lemma:convexity_growth}
  Let $g:\R^d\to\R$ be a convex function. Fix $\theta_0,\theta_1\in\R^d$ and
  $t\geq 1$, and define $\theta_t = \theta_0 + t(\theta_1 - \theta_0)$. Then
  $g(\theta_t) \geq g(\theta_0) + t(g(\theta_1) - g(\theta_0))$.
\end{lemma}
\begin{proof}
  Rearranging terms, we have $\theta_1 = \frac 1 t \theta_t + (1 - \frac 1
  t)\theta_0$ . Applying the definition of convexity at $\theta_1$ gives
  \[g(\theta_1) \leq \frac 1 t g(\theta_t) + \left(1-\frac 1 t \right)g(\theta_0),\]
  which is equivalent to $g(\theta_t) \geq g(\theta_0) + t(g(\theta_1) - g(\theta_0))$.
\end{proof}

Returning to Proposition~\ref{prop:consistency-convex}, the local
Lipschitzness condition implies there is an $\epsilon > 0$ such that
\begin{align}
  \label{eq:convergence_epsball}
  \sup_{\theta:\|\theta-\theta\opt\|\leq \epsilon} |\LPP_{\lambda}(\theta) - L(\theta)| \cp 0
\end{align}
by a standard covering argument.
%
By the uniqueness of $\theta\opt$, for all
$\epsilon>0$ we know there exists a $\delta>0$ such that $L(\theta) -
L(\theta\opt)\geq \delta$ for all $\theta$ on the $\epsilon$-shell $\{\theta
\mid \norm{\theta -
\theta\opt}=\epsilon\}$ around $\theta\opt$.
%
With this we can write
\begin{align*}
  \inf_{\norm{\theta-\theta\opt} = \epsilon}
  \LPP_{\lambda}(\theta) - \LPP_{\lambda}(\theta\opt)
  & = \inf_{\norm{\theta - \theta\opt} =\epsilon} \left((\LPP_{\lambda}(\theta) - L(\theta) + (L(\theta) - L(\theta\opt)) + (L(\theta\opt) - \LPP_{\lambda}(\theta\opt))\right)\\
  &\geq \delta - o_P(1),
\end{align*}
where the convergence~\eqref{eq:convergence_epsball}
guarantees that the first and third terms above vanish.

Now fix any $\theta$ such that $\norm{\theta-\theta\opt} \geq \epsilon$. We
apply Lemma~\ref{lemma:convexity_growth} with $\theta_0 = \theta\opt$ and
$\theta_1 = \frac{\theta-\theta\opt}{\|\theta-\theta\opt\|}\epsilon$.
Combining the preceding
display with Lemma~\ref{lemma:convexity_growth} gives
\begin{equation*}
  \LPP_{\lambda}(\theta) - \LPP_{\lambda}(\theta\opt) \geq \frac{\|\theta-\theta\opt\|}{\epsilon}(\delta-o_P(1)) \geq \delta - o_P(1).
\end{equation*}
Therefore, no $\theta$ with $\|\theta-\theta\opt\|\geq \epsilon$ can
minimize $\LPP_{\lambda}$, and so by contradiction we have shown
$\P(\norms{\thetaPP_{\lambda} - \theta\opt}< \epsilon)\to 1$, as desired.

We turn to the final claim of the proposition, which allows random
$\hat{\lambda}$. The convergence~\eqref{eq:convergence_epsball} holds
uniformly for $\lambda'$ in a small neighborhood of $\lambda$ as well
because of an identical covering argument; as $\hat{\lambda} = \lambda +
o_P(1)$, with probability tending to 1 we have that $\hat{\lambda}$ belongs to
the neighborhood of $\lambda$. The rest of the proof is identical.


\subsection{Proof of Theorem \ref{theorem:confidence-set-equivalence}}
\label{sec:proof-confidence-set-equivalence}

We will require that the loss function is stochastically smooth in the following sense.

\begin{assumption}[Stochastic smoothness]
  \label{assumption:donsker}
  There exists a compact neighborhood $B$ of $\theta\opt$ such that
  the classes
  \begin{equation*}
    \mc{L} \defeq \{(x, y) \mapsto \nabla \loss_\theta(x, y)
    : \theta \in B \}
    ~~~ \mbox{and} ~~~
    \mc{L}^f
    \defeq \{x \mapsto \nabla \loss_\theta(x, f(x)) : \theta \in B\}
  \end{equation*}
  are both $\P$-Donsker.
\end{assumption}



Let $Z \sim \normal(0, I)$. We will show that,
  for any $c < \infty$,
  \begin{equation}
  \label{eq:thm2_proofgoal}
    \limsup_n
    \sup_{\norm{h} \le c}
    \left|
    \P(\theta\opt + h/\sqrt{n} \in \CPP_\alpha)
    - \P\big(\ltwobig{-\Sigma^{-1/2} h + Z}^2
    \le \chi^2_{d,1-\alpha} \big) \right|
    = 0,
  \end{equation}
  for $\CPP_\alpha \in\{ \C_\alpha^{\mathrm{PP},T}, \C_\alpha^{\mathrm{PP},U}\}$, where $\Sigma$ is the asymptotic covariance from Theorem \ref{theorem:normality} for $\lambda = 1$, i.e. $\Sigma \equiv \Sigma^1 = H_{\theta\opt}^{-1} (r V^1_{f,\theta\opt} + V^1_{\Delta,\theta\opt}) H_{\theta\opt}^{-1}$.
  
The result \eqref{eq:thm2_proofgoal} directly implies the main result of the theorem.
  
  
  
  
  
  Let $h_n$ be any sequence such that $\sqrt{n} h_n \to h$ for some fixed $h
  \in \R^d$.  We give alternative representations of the statistics $T_n$
  and $U_n$ under the local alternative parameters $\theta\opt + h_n$. We
  leverage the following two technical lemmas:
  
  
  
 \begin{lemma}
    \label{lemma:variance-control}
    Under the conditions of Theorem~\ref{theorem:confidence-set-equivalence},
    \begin{equation*}
      \VhatDelta(\theta\opt + h_n)
      \cp \cov(\nabla \loss_{\theta\opt}
      - \nabla \loss_{\theta\opt}^f)
      ~~ \mbox{and} ~~
      \Vhatf(\theta\opt + h_n)
      \cp \cov(\nabla \loss_{\theta\opt}^f).
    \end{equation*}
  \end{lemma}
  
 
 

\begin{lemma}
    \label{lemma:gradient-process-linear}
   Define the (sub)gradient process
\begin{equation}
  \label{eqn:gradient-process}
\widehat G(\theta) \defeq \E_n[\nabla \loss_{\theta}
  - \nabla \loss_{\theta}^f] - \E[\nabla \loss_{\theta}
  - \nabla \loss_{\theta}^f] + \widetilde \E_N[\nabla \loss_{\theta}^f] - \E[\nabla \loss_{\theta}^f].
 \end{equation}
Then, under the conditions of Theorem~\ref{theorem:confidence-set-equivalence}, $\widehat G(\theta\opt + h_n) = \widehat G(\theta\opt) + o_P(1/\sqrt{n}).$
\end{lemma}

Lemma \ref{lemma:variance-control} follows by the classical
  result~\cite[Lemma 2.10.14]{VanDerVaartWe96} that if $\mc{F}$ is a Donsker
  class, then $\mc{F}^2 = \{f^2 : f \in \mc{F}\}$ is Glivenko-Cantelli.
  
  Lemma~\ref{lemma:gradient-process-linear} follows by the fact that the process is asymptotically stochastically equicontinuous, which follows by Assumption \ref{assumption:donsker}; see below for the proof.
  
\begin{proof}[Proof of Lemma \ref{lemma:gradient-process-linear}]
We adopt the definitions of $\GG_n g$, $\GG_N^f g$, and $\GG_n^f g$ from Equations \eqref{eq:Gn}, \eqref{eq:GNf}, and \eqref{eq:Gnf}, respectively.
    Because
    the process $[\GG_n \nabla \loss_\theta]_{\theta \in B}$
    converges to a (tight) Gaussian process in $\theta$
    by Assumption~\ref{assumption:donsker},
    as $h_n = O_P(1/\sqrt{n})$ we necessarily have
    \begin{equation*}
      \GG_n \nabla \loss_{\theta\opt + h_n}
      - \GG_n\nabla \loss_{\theta\opt}
      \cp 0
    \end{equation*}
    (as the process is asymptotically stochastically equicontinuous
    by Assumption~\ref{assumption:donsker}). Similarly,
    Assumption~\ref{assumption:donsker} implies $\GG_N^f \nabla \loss_{\theta\opt + h_n}
      - \GG_N^f\nabla \loss_{\theta\opt} \cp 0$ and $\GG_n^f \nabla \loss_{\theta\opt + h_n}
      - \GG_n^f\nabla \loss_{\theta\opt} \cp 0$.
    Adding and subtracting the preceding quantities and rescaling by $\sqrt{n}$,
    we obtain the lemma.
\end{proof}  
  
  
  
  
%  Now we use Lemma~\ref{lemma:gradient-process-linear} to give the
%  asymptotically equivalent formulations of $T$ and $U$
%  evaluated at the alternatives $\theta\opt + h_n$.
%By the analysis in Theorem \ref{theorem:normality}, we have
%  \begin{equation*}
%    \thetaPP = \theta\opt
%    + \nabla^2 L(\theta\opt)^{-1}
%    \widehat G(\theta\opt)
%    + o_P(1/\sqrt{n}).
%  \end{equation*}
Now we give asymptotically equivalent formulations of $T_n$ and $U_n$ evaluated at the alternatives $\theta\opt + h_n$. We can write $T_n$ as
\begin{equation*}
  T_n(\theta\opt + h_n)
  = \sqrt{n} \widehat\Sigma^{-1/2} (\thetaPP - \theta\opt - h_n)
  = \sqrt{n} \widehat\Sigma^{-1/2} (\thetaPP - \theta\opt)
  - \Sigma^{-1/2} h + o_P(1).
\end{equation*}
Turning to $U_n$, we let $V_{\Delta,\theta\opt} = \cov(\nabla \loss_{\theta\opt}
- \nabla \loss_{\theta\opt}^f)$ and
$V_{f,\theta\opt} = \cov(\nabla \loss_{\theta\opt}^f)$. Then Lemma \ref{lemma:variance-control} and
Lemma~\ref{lemma:gradient-process-linear} give
\begin{align*}
  &U_n(\theta\opt + h_n)\\
  &\quad \stackrel{(i)}{=} \sqrt{n} (V_{\Delta,\theta\opt} + \ratio V_{f,\theta\opt} + o_P(1))^{-1/2}
  \widehat G(\theta\opt + h_n)
  + \sqrt{n} (V_{\Delta,\theta\opt} + \ratio V_{f,\theta\opt} + o_P(1))^{-1/2}
  \nabla L(\theta\opt + h_n)
  \\
  &\quad \stackrel{(ii)}{=}
  \sqrt{n} (V_{\Delta,\theta\opt} + \ratio V_{f,\theta\opt})^{-1/2}
  \widehat G(\theta\opt) + (V_{\Delta,\theta\opt} + \ratio V_{f,\theta\opt})^{-1/2} H_{\theta\opt} h
  + o_P(1),
\end{align*}
where step~$(i)$ follows from Lemma~\ref{lemma:variance-control}
and $(ii)$ from Lemma~\ref{lemma:gradient-process-linear} coupled
with Slutsky's lemmas and
that $\nabla L(\theta\opt + h_n) = \nabla L(\theta\opt + h_n)
- \nabla L(\theta\opt) = \nabla^2 L(\theta\opt)h_n + o(\norm{h_n})
= H_{\theta\opt} h_n + o(1/\sqrt{n})$.

Rewriting the preceding two displays and applying Theorem \ref{theorem:normality} to the first,
we have
\begin{equation*}
  T_n(\theta\opt + h_n) \cd \normal(-\Sigma^{-1/2} h, I)
  ~~ \mbox{and} ~~
  U_n(\theta\opt + h_n) \cd \normal((V_{\Delta,\theta\opt} + \ratio V_{f,\theta\opt})^{-1/2} H_{\theta\opt} h, I).
\end{equation*}
The squared norm of the latter mean satisfies
\begin{equation*}
  \ltwo{(V_{\Delta,\theta\opt} + \ratio V_{f,\theta\opt})^{-1/2} H_{\theta\opt} h}^2
  = h^T H_{\theta\opt} (V_{\Delta,\theta\opt} + \ratio V_{f,\theta\opt})^{-1}
  H_{\theta\opt} h =
  h^T \Sigma^{-1} h
\end{equation*}
by the definition of $\Sigma$. Therefore, the rotational invariance of the Gaussian yields
\begin{equation}
  \begin{split}
    \P(\theta\opt + h_n \in \C_\alpha^{\mathrm{PP},T})
    & = \P(\ltwo{T_n(\theta\opt + h_n)}^2
    \le \chi^2_{d,1-\alpha})
    = \P\left(\ltwo{-\Sigma^{-1/2} h + Z}^2 \le
    \chi^2_{d,1-\alpha}\right) + o(1) \\
    \P(\theta\opt + h_n \in \C_\alpha^{\mathrm{PP},U})
    & = \P(\ltwo{U_n(\theta\opt + h_n)}^2
    \le \chi^2_{d,1-\alpha})
    = \P\left(\ltwo{-\Sigma^{-1/2} h + Z}^2 \le
    \chi^2_{d,1-\alpha}\right) + o(1),
  \end{split}
  \label{eqn:almost-home}
\end{equation}
where $Z \sim \normal(0, I)$.

Finally, we employ a standard compactification argument. Assume for the sake
of contradiction that
\begin{equation*}
  \limsup_n \sup_{\norm{h} \le c}
  \left|\P(\theta\opt + h / \sqrt{n} \in \CPP_\alpha)
  - \P\left(\ltwo{-\Sigma^{-1/2} h + Z}^2 \le \chi^2_{d,1-\alpha}\right)
  \right| > 0,
\end{equation*}
for one of $\CPP_\alpha \in \{\C_\alpha^{\mathrm{PP},T}, \C_\alpha^{\mathrm{PP},U}\}$.
If this is the case, then there must be a bounded subsequence of $h_n$
with
\begin{equation*}
  \left|\P(\theta\opt + h_n / \sqrt{n} \in \CPP_\alpha)
  - \P\left(\ltwo{-\Sigma^{-1/2} h_n + Z}^2 \le \chi^2_{d,1-\alpha}\right)
  \right| \to a > 0.
\end{equation*}
Take a convergent subsequence from this set, so that $\sqrt{n} (h_n / \sqrt{n})
= h_n \to h$. The existence of this subsequence contradicts
the limits~\eqref{eqn:almost-home}. This proves the first claim of the theorem.


For the second claim, take any $t_n = \omega(1/\sqrt{n})$ and $\delta\in(0,1)$. Then
\begin{align*}
  \P\left( \theta\opt + h t_n \in \C_\alpha^{\mathrm{PP},T}\right) &\leq \P\left( \|\sqrt{n}\widehat\Sigma^{-1/2}  h t_n \| \leq \chi_{d,1-\alpha} + \|\sqrt{n}\widehat\Sigma^{-1/2}  (\thetaPP - \theta\opt) \|\right)\\
  &\leq \P\left( \|\sqrt{n}\widehat\Sigma^{-1/2}  h t_n \| \leq \chi_{d,1-\alpha} + \chi_{d,1-\delta}\right) + \P\left( \|\sqrt{n}\widehat\Sigma^{-1/2}  (\thetaPP - \theta\opt) \| > \chi_{d,1-\delta}\right).
\end{align*}
Taking the limit, we get
\begin{equation*}
  \limsup_n \sup_{\|h\|\geq c} \P\left( \theta\opt + h t_n \in \C_\alpha^{\mathrm{PP},T}\right) \leq  \delta.
\end{equation*}
As $\delta$ was arbitrary, we conclude that $\limsup_n \sup_{\|h\|\geq c}
\P( \theta\opt + h t_n \in \C_\alpha^{\mathrm{PP},T}) = 0$.



%% \begin{align*}
%%   0 & = \E_t[\nabla \loss_{\theta_t}]
%%   + \lambda \big(\wt{\E}_t^f[\nabla \loss_{\theta_t}]
%%   - \E_t^f[\nabla \loss_{\theta_t}]\big) \\
%%   & = \E_0[(1 + t g) \nabla \loss_{\theta_t}]
%%   + \lambda \big(\wt{\E}^f_0[(1 + t g)\nabla \loss_{\theta_t}]
%%   - \E_0^f[(1 + t g)\nabla \loss_{\theta_t}]\big) \\
%%   & \stackrel{(\star)}{=} \left(\E_0[\nabla^2 \loss_{\theta_0}] + o(1)\right)
%%   (\theta_t - \theta_0)
%%   + \E_0[(1 + t g) \nabla \loss_{\theta_0}] \\
%%   & \qquad ~ + t \lambda \left(
%%   \wt{\E}_0[g(X, Y) \nabla \loss_{\theta_t}(X, f(X))]
%%   - \E_0[g(X, Y) \nabla \loss_{\theta_t}(X, f(X))]\right)
%% \end{align*}
%% as $t \to 0$, where equality~$(\star)$ uses the smoothess of the losses and
%% that $\wt{\E}_0[\nabla \loss_\theta(X, f(X))] - \E_0[\nabla \loss_\theta(X,
%%   f(X))] = 0$. Once again using that $\nabla \loss_\theta$ is smooth in
%% $\theta$, from the optimality condition
%% $\E_0[\nabla\loss_{\theta_0}] = 0$ and one more
%% Taylor approximation we deduce
%% \begin{equation*}
%%   0 = \left(\nabla^2 L(\theta_0) + o(1)\right)
%%   (\theta_t - \theta_0)
%%   + t \E_0[g \nabla \loss_{\theta_0}]
%%   + t \lambda\left(\wt{\E}_0 - \E_0\right)[g(X, Y)
%%     \nabla \loss_{\theta_0}(X, f(X))]
%%   - o(t),
%% \end{equation*}
%% and rearranging,
%% \begin{equation*}
%%   \theta_t - \theta_0
%%   = t \cdot \nabla^2 L(\theta_0)^{-1}
%%   \left(
%%   \E_0\left[g(X, Y) (\nabla \loss_{\theta_0}
%%     - \lambda \nabla \loss_{\theta_0}^f)\right]
%%   + \lambda \wt{\E}_0[g(X, Y) \nabla \loss_{\theta_0}^f)]\right)
%%   + o(t).
%% \end{equation*}
%% In particular, let $a \in \{0, 1\}$ indicate whether a data point
%% $(x, y)$ comes from the labeled distribution $P_0$ ($a = 0$) or
%% the unlabeled data $\wt{P}_0$ (so that abusing notation,
%% $(x, y) = (x, f(x))$ in this case).
%% Then for fixed $\lambda$, the parameter
%% $\theta$ has influence function
%% \begin{equation*}
%%   \dot{\theta}(x, y, a) \defeq
%%   \begin{cases}
%%     \nabla^2 L(\theta_0)^{-1}
%%     \left(\nabla \loss_{\theta_0}(x, y) - \lambda (\nabla \loss_{\theta_0}(x, f(x))
%%     - \E^f[\nabla \loss_{\theta_0}]
%%     \right) & \mbox{if~} a = 0 \\
%%     \nabla^2 L(\theta_0)^{-1}
%%     \lambda \left(\nabla \loss_{\theta_0}(x, f(x)) - \E^f[\nabla \loss_{\theta_0}]
%%     \right) & \mbox{if~} a = 1.
%%   \end{cases}
%% \end{equation*}

%% Substituting back for $\theta_0 = \theta\opt$,
%% this has asymptotic covariance as $n, N \to \infty$,
%% \begin{equation*}
%%   \Sigma^\lambda
%%   = H_{\theta\opt}^{-1}
%%   \left(r \lambda^2 \var(\nabla \loss^f_{\theta\opt})
%%   + \var(\nabla \loss_{\theta\opt} - \lambda \nabla \loss_{\theta\opt}^f)
%%   \right) H_{\theta\opt}^{-1}.
%% \end{equation*}

%% \jcdcomment{Here, write out the influence function stuff... TODO}



\subsection{Proof of Proposition \ref{prop:optimal-lambda}}
\label{sec:proof-optimal-lambda}

By definition
\begin{equation*}
  \lambda\opt = \argmin_{\lambda} \Tr\left(\Sigma^\lambda\right) = \argmin_{\lambda} \Tr\left(H_{\theta\opt}^{-1} \left(r \cdot \lambda^2 \cov(\nabla \loss_{\theta\opt}^f) + \cov( \nabla \loss_{\theta\opt} -  \lambda \nabla \loss_{\theta\opt}^f ) \right) H_{\theta\opt}^{-1}\right).
\end{equation*}
Writing $\cov( \nabla \loss_{\theta\opt} - \lambda \nabla
\loss_{\theta\opt}^f ) = \cov(\nabla \loss_{\theta\opt}) + \lambda^2
\cov(\nabla \loss_{\theta\opt}^f) - \lambda (\Cov(\nabla
\loss_{\theta\opt}, \nabla \loss_{\theta\opt}^f) + \Cov(\nabla
\loss_{\theta\opt}^f, \nabla \loss_{\theta\opt} ))$ and applying the
linearity of the trace, $\lambda\opt$ evidently minimizes
\begin{align*}
\lambda^2(1+r) \Tr\left(H_{\theta\opt}^{-1} \cov(\nabla \loss_{\theta\opt}^f)H_{\theta\opt}^{-1}\right) - \lambda \Tr\left(H_{\theta\opt}^{-1}\left(\Cov\left(\nabla \loss_{\theta\opt}, \nabla \loss_{\theta\opt}^f \right) + \Cov\left(\nabla \loss_{\theta\opt}^f, \nabla \loss_{\theta\opt} \right) \right)H_{\theta\opt}^{-1} \right).
\end{align*}
Analytically optimizing this quadratic gives the optimal choice
\begin{equation*}
  \lambda\opt = \frac{\Tr\left(H_{\theta\opt}^{-1}\left(\Cov\left(\nabla \loss_{\theta\opt}, \nabla \loss_{\theta\opt}^f \right) + \Cov\left(\nabla \loss_{\theta\opt}^f, \nabla \loss_{\theta\opt} \right) \right)H_{\theta\opt}^{-1} \right)}{2(1+r) \Tr\left(H_{\theta\opt}^{-1}\cov(\nabla \loss_{\theta\opt}^f) H_{\theta\opt}^{-1} \right) }.
\end{equation*}
